\section{Introduction}\label{sec:r23intro}
A numerical solution of a dynamic economic model has to deliver a policy, improve its economic payoff, and control its loss relative to the original optimum. A neural representation can change the geometry of policy generation without changing the attainable policies. A value witness can sharpen a comparison without improving the policy it accompanies. Confusing these effects obscures both the economic accuracy of a solver and the source of a computational advantage.

We develop Neural Bellman Operators that distinguish these effects while retaining a common stochastic control problem. The application is a stopped consumption--portfolio economy with costly preference adjustment, correlated shocks, continuous controls, and a time-dependent liquidation settlement. Its unrestricted objective, full state--time domain, and accuracy target of $.01$ are unchanged. The operator couples a candidate-generation map to policy-specific payoff acceptance and a separate unrestricted value comparison. The former determines whether a proposed policy should replace its incumbent; the latter determines the policy's certified economic error.

Our first contribution is a precise analysis of the financed policy coordinates. The budget equation removes one common-intercept direction. We establish its unique solution, its implicit derivative, finite price-root brackets that include unbounded Gaussian tails, and an explicit inverse-conditioning bound. The derivative is an oblique projection, not an arbitrary normalization. We also prove that the neural and direct parameterizations can represent the same finite-node raw-output arrays. Consequently, a comparison between them identifies a difference in parameterization and optimization geometry, not an unmeasured difference in the size of the delivered policy class. A directed full-rank certificate verifies a concrete interpolation construction.

Our second contribution is a continuum payoff comparison for independently initialized expert policies. Verified pair moments control the old mixture's Jensen benefit when both its consumption slopes and preference drifts differ across corners. The resulting theorem proves policy improvement and a true-regret reduction at every initial state of a rectangle. Correlated shocks do not require an independence approximation: their contribution is bounded on the joint time--probability measure. The unrestricted upper comparison is used for regret but is unnecessary for the mathematical ordering of two feasible policy payoffs.

The experiments separate parameterization from optimizer and include a control that removes the redundant budget coordinate. The new held-out design contains 72 configurations: three independently initialized four-expert systems, three coordinate representations, and two optimizers. Each configuration has the same 400-gradient-call cap. The four primary neural/direct arms start from exactly the same delivered policy within each block. The 47-coordinate direct control is mapped to that policy, with the numerical discrepancy recorded and its incumbent separately checked whenever it is not identical. All configurations and all prescribed comparisons are reported.

In this design every neural--Adam system has regional regret below \RXXIIINeuralRegret{} and a uniform initial-to-final payoff gain of at least \RXXIIIGainMin{}. Its uniform payoff advantage over each of the two direct--Adam coordinate systems is at least \RXXIIIConditionalGain{}. Removing the redundant intercept does not remove this conditional advantage. Direct L-BFGS-B nevertheless reaches a tighter regional bound with substantially fewer objective calls, and one neural L-BFGS-B system remains above $.01$. The evidence identifies a reproducible optimizer-conditional effect; it does not establish an unconditional neural accuracy--cost frontier advantage.

The full-state experiment is a distinct numerical object. Its actor and witness take the current $(t,u,x)$ as inputs and are checked on every cell of the original domain against the original continuous action set. A recorded warm-start refinement tightens its all-start-time certificate from $7.2783190635$ to \RXXIIGlobalBound{}. This is an improvement in a certified comparison, not a separately proved increase in that actor's payoff. We retain the full-domain $.01$ objective rather than substitute the successful regional experiment for it. A constructive residual-correction cone, valid without exact financial replication, makes the requirements for further global improvement explicit.

The paper also separates mathematical deployment from numerical implementation. Regional policies are compiled to slab coefficients; their conditional prices and hedges still require integration. The selected consumption streams are replicable, but the preference shock has an unspanned component: exact replication of those streams does not make every payoff in the joint filtration replicable. The independent dual upper is non-neural and its construction cost is charged. A wealth counterfactual is a sufficient re-budgeted increment for fixed learned shapes, not a minimum compensating variation for an unchanged policy map.

The numerical dynamic-programming and verification background includes \citet{Judd1998}, \citet{KushnerDupuis2001}, and \citet{BarlesSouganidis1991}; adaptive sparse-grid comparisons are developed by \citet{BrummScheidegger2017}. Dual comparisons relate to \citet{BrownSmithSun2010,BrownSmith2011}; neural control and high-dimensional equation methods include \citet{HanE2016,HanJentzenE2018}. Our results combine financed-coordinate identification, nonsynchronized continuum comparison, and executed acceptance with an explicit accounting of their distinct roles. Earlier theory, numerical failures, and high-dimensional comparisons remain in the appendices and supplement with their original scopes. They are not relabeled as evidence from the new held-out study.
